Two frameworks for dealing with uncertainty in time intervals have been presented. In this section, we will study the equivalences between an uncertain time interval (UTI) and an ill-known time interval (IKTI). Nevertheless, there are some differences in the semantics of UTIs and IKTIs which are also explained. Next to that, the computation of the Allen's Relations are compared.

\subsection{\label{subsec:sim-dif}Similarities and differences}
A time interval is defined by a starting and an ending point. When uncertainty arises, in one or both starting and / or ending points, then either an UTI $\pi_I$ or an IKTI can be used. 

\begin{itemize}
 \item An IKTI $J = \left[X, Y \right]$ is defined by two ill-known values $X, Y$ which are respectively, the starting and ending points of the time interval. Two possibility distributions $\pi_X, \pi_Y$ are defined for $X$ and $Y$ respectively.
\item An UTI $\pi_I$ is defined by two possibility distributions, $\pi_S$, $\pi_E$ for the starting  and ending point respectively.
\end{itemize}

In the following, we will assume that:

\begin{align}
\pi_X &= \pi_S = \mu_S \\
\pi_Y &= \pi_E = \mu_E
\end{align}

Where $\mu_S, \mu_E$ are triangular membership functions. Notice that the main difference between an UTI and an IKTI is a matter of interpretation. IKTIs  keep both possibility distributions for both starting and ending poins, while UTIs take a t-norm (the minimum) of both possibility distributions (see Eq. \eqref{eq:possibility-uti}). That is, not all the crisp intervals between the boundaries of the starting and ending point are candidate intervals for the UTI.

\begin{figure}[h]
   \centering
   \includegraphics[scale=1.5]{graphs/equivalences.eps}
   \caption{Possibility degrees for two candidate intervals with respect to the UTI $\pi_I$. The upper crisp time interval has a possibility degree of 1, whereas the lower crisp interval has a possibility degree of 0. }
   \label{fig:ikc-greater}
 \end{figure}

It is possible to obtain the same interpretation for an IKTI, but some ill-known constraints must be defined over the starting and ending ill-known values $X$, $Y$. In \cite{Pons2011}, two constraints are defined to determinate whether a crisp time interval $[a, b]$ resides between the boundaries of $J$. In this case, all the crisp intervals whose boundaries are within the boundaries of $X$ and $Y$ are considered to be a candidate interval. Nevertheless, if we want to obtain the same interpretation as in the UTI, the following two ill-known constraints must be defined:

\begin{align}
 \label{eq:ikc-ikti-uti}
C_1 &\triangleq \left(=, X \right) \\
C_2 &\triangleq \left(=, Y \right)\\
\Pos \left(\left[a, b\right] \right) &= C_1 \left(a \right) \wedge C_2\left(b \right) =\\
\label{eq:equivalence-ikti-uti}
\Pos \left(\left[a, b\right] \right) &= \min \left( \pi_X(a), \pi_Y(b) \right)
\end{align}

In this case, Eq. \eqref{eq:equivalence-ikti-uti} is equivalent to Eq. \eqref{eq:possibility-uti}.
In the following subsection, we will study the connection between the Allen's relations in each framework.

\subsection{\label{subsec:eval-allen-relations}Evaluation of the Allen's Relations}
As it has been shown in Table \ref{tab:urz}, when using UTIs in the triangular model, it is possible to obtain a more detailed relational information (there are fifteen possible relations) than with the ill-known constraint approach. The following observations should be taken into account when extracting relational information from both frameworks:

\begin{itemize}
 \item In TM, the Allen's Relations $R \in \left \lbrace B, O, C, D, OB, A \right \rbrace$ between an interval point $J$  and an UTI $\pi_I$ are computed by Eq. \eqref{eq:possibility-iz-uti-certain}. Therefore, for the Allen's Relations in $R$, 
 possibility that an Allen Relation between an interval point $J$  and an UTI $\pi_I$ is $1$ if the interval point $J$ is in the Allen Relation, or $0$ in other case.
\item In the TM, the other nine Allen's Relations between an interval point $J$ and an UTI $\pi_I$ are graded from $1$ to $0$ and computed by using the Eq. \eqref{eq:possibility-iz-uti}. 
\item In the ill-known constraint framework, all the thirteen Allen's Relations are computed by using an aggregation of ill-known constraints (see Table \ref{tab:allen-relations}). Therefore, the possibility that a time point $J$ is in a given Allen Relation is always graded from $1$ to $0$.
\end{itemize}

For example, consider that we want to know whether the relation Before (B) holds between a time point $J$ and an IKTI $I$. We will note $\Pos \left(C_B (J)  \right)$ the possibility that the relation before holds. There are three possibilities:
\begin{enumerate}
 \item $\Pos \left(C_B (J)  \right)$ = 1: The time point $J$ is before the IKTI $I$.
\item $\Pos \left(C_B (J)  \right)$ = 0: The time point $J$ is not before the IKTI $I$.
\item $0 < \Pos \left(C_B (J)  \right) < 1$: The time point $J$ is to some extent before the IKTI $I$.
\end{enumerate}

The last case corresponds in the TM to case when the before relation is one of the possible relations (See Table \ref{tab:urz}). That is, one of $\left \lbrace PM, PMB, UB \right \rbrace$.

The following table establishes the correspondences between the Allen's relations in the TM and in the ill-known constraints framework.




\begin{table}[h]
\centering
\begin{tabular}{|l|l|l|}
\hline
Symbol & Name & Possible Relations \\
\hline
B & Before & B, PM, PMB, UB \\
\multirow{2}{*}
{O} & Overlaps & O, PM, PS, PFB,\\
     &          &  PE, PMB, UB \\
C & Contains & C, PFB, PE, PSB \\
\multirow{2}{*}
D & During & D, PS, PE, \\
  &         &  PMB, UB  \\
OB & Overlapped by & OB, PE, PSB, PF \\
A & After & A, PF, UE \\
M & Meets & M, PM, PMB, UB \\
MB & Met By & MB, PF, UE \\
E & Equals & E, PE \\
\multirow{2}{*}
S & Starts & S, PS, PE, \\
  &         &   PMB, UB \\
SB & Started by & SB, PE, PSB \\
F & Finishes & F, PE \\
FB & Finished By & FB, PE, PFB\\
\hline
\end{tabular}
\caption{Correspondences between the Allen's Relations in the ill-known constraint framework and the triangular model,}
\label{tab:correspondences-tm-ikcf}
\end{table}

 
We will illustrate the calculation of one Allen Relation with an example.
\begin{example}
Consider a crisp time interval $J = [1, 3]$ the two ill-known time points $S, E$ that define a time interval. 
We want to obtain the possibility that $J$ is before the time interval.
The points are defined by two triangular membership functions $\mu_S$, $\mu_E$.

\begin{align}
\label{eq:ex-point}
\pi_S &= \mu_S = \left[2,5,7 \right] \\
\pi_E &= \mu_E = \left[9,12,15 \right]
\end{align}

Then, it is possible to define both an IKTI $I$ Eq. \eqref{eq:ex-uti-ikti} and an UTI $\pi_I$ Eq. \eqref{eq:ex-uti-uti}. The intervals are illustrated in Figure \ref{fig:example}.

\begin{figure}[h]
   \centering
   \includegraphics[scale=1.5]{graphs/example.eps}
   \caption{In the top, the crisp time interval $J$ and the IKTI given by $I = [S, E]$. The ill-known constraint for the before relation is illustrated in grey color. In the lower graphic, the interval point $J$ and the UTI $\pi_I$. }
   \label{fig:example}
 \end{figure}

\begin{align}
\label{eq:ex-uti-ikti}
I = \left[S, E \right]\\
\label{eq:ex-uti-uti}
\mu_I = \left(\mu_S, \mu_E \right)
\end{align}

First we will obtain the possibility between $J$ and the UTI $\pi_I$. The relation possibly meets has three candidates: Before, Meets and Overlaps.

\begin{align}
  \mu_{\left(J R \pi_I \right)}(B) &= \sup_{a < e} \pi_S(a)\\
\mu_{\left(J R \pi_I \right)}(M) &= \pi_S(e)\\
\mu_{\left(J R \pi_I \right)}(O) &= \sup_{a > e} \pi_S(a)
\end{align}

\begin{align}
\mu_{\left(J R \pi_I \right)}(B) &= 0.67\\
\mu_{\left(J R \pi_I \right)}(M) &= 0.33\\
\mu_{\left(J R \pi_I \right)}(O) &= 0
\end{align}

In order to compute the possibility  for the before relation between the crisp interval $J$ and the IKTI $I$, a ill-known constraint is needed.

\begin{equation}
 C_1\stackrel{\triangle}{=} \left(<,X\right)
\end{equation}

Then the possibility is computed as follows:

\begin{align}
\Pos \left(C_1 (J) \right) &= \sup_{a < e} \pi_S(a)\\
\Pos \left(C_1 (J) \right) &= 0.67 
\end{align}






\end{example}



% 
% In this section we will analyze first the calculation of all the possible Allen's relations between a crisp time point and an ill-known time point and between a crisp time point and an uncertain time indication (UTI). Then, we will compare two ill-known time points and two uncertain time indications.

% \subsection{\label{subsec:prop-crisp-vs-ill}Calculation of crisp time intervals}
% In this section we first study how uncertainty about wheter crisp interval is enclosed by the ill-known interval given by $[X,Y]$ or by the UTI.
% 
% \subsubsection{Crisp time interval enclosed by an ill-known time interval}
% Consider a crisp time interval given by $I=[a, b]$. We want to know whether all the points in this interval are within the boundaries of the ill-known interval given by $[X, Y]$. The calculation is done by means of two ill-known constraints as explained in \cite{}:
% \begin{eqnarray}
% C_1\stackrel{\triangle}{=}\left(\geq,X\right)\\
% C_2\stackrel{\triangle}{=}\left(\leq,Y\right).
% \end{eqnarray}
% 
% The possibility and the necessity are given by:
% \begin{eqnarray}
% \Pos\left(\lambda([a,b])\right)&=&\min\bigg(\Pos(C_1([a,b])),\Pos(C_2([a,b]))\bigg)\\
% \Nec\left(\lambda([a,b])\right)&=&\min\bigg(\Nec(C_1([a,b])),\Nec(C_2([a,b]))\bigg).
% \end{eqnarray}
% Which can be expanded to:
% \begin{eqnarray}
% \label{eq:interval-pos}
% \Pos\left(\lambda([a,b])\right)&=&\min\bigg(\sup_{a\leq w}\pi_{X}(w),\sup_{b\geq w}\pi_{Y}(w)\bigg)\\
% \label{eq:interval-nec}
% \Nec\left(\lambda([a,b])\right)&=&\min\bigg(\inf_{a>w}1-\pi_{X}(w),\inf_{b<w}1-\pi_{Y}(w)\bigg).
% \end{eqnarray}
% 
% 
% \subsubsection{\label{subsubsec:crisp-uti}Crisp time interval enclosed by an UTI}
% Consider now a crisp time interval given by $I=[a, b]$ and an Uncertain Time Interval given by $\pi_I = \left(\pi_s, \pi_e \right)$. Again, we want to know if all the points in the interval are within the boundaries of the UTI. The calculation is done as explained in \cite{}:
% \begin{equation}
%  \pi_I \left( \left[a, b  \right] \right) =
% \begin{cases}
%  \min \left(\pi_s(a),\pi_e(b)\right) & \mbox{ if } a \leq b\\
% 0 				     & \mbox{ otherwise. }
% \end{cases}
% \end{equation}

